\documentclass[11pt]{article}
\usepackage{preamble}

\newcommand{\lessonrow}[2]{\textbf{#1} & #2 \\ \hline}

\begin{document}

\packetheading{Lesson 5-4 \textperiodcentered\ Period 1 \textperiodcentered\ Student Packet}{Solving Radical Equations --- Foundation}

\sectionbanner{OBJECTIVES}
{\small
\renewcommand{\arraystretch}{1.25}
\noindent\begin{tabular}{|p{1.8in}|p{4.8in}|}
\hline
\lessonrow{Math Objective}{Solve radical equations by isolating the radical and raising both sides to the appropriate power; identify and reject extraneous solutions.}
\lessonrow{Language Objective}{Explain in writing, using sentence frames, why squaring both sides can introduce extraneous solutions and how to check.}
\lessonrow{Essential Understanding}{Squaring both sides is NOT a reversible operation --- it can create solutions that don't satisfy the original equation. Every candidate must be checked.}
\end{tabular}
}

\sectionbanner{TOPIC GOALS / ESSENTIAL QUESTION / MATERIALS}
{\small
\renewcommand{\arraystretch}{1.25}
\noindent\begin{tabular}{|p{1.8in}|p{4.8in}|}
\hline
\lessonrow{Topic Goals}{Solve single-radical equations; detect extraneous solutions.}
\lessonrow{Essential Question}{Why does squaring both sides sometimes produce solutions that don't work, and how do we catch them?}
\lessonrow{Materials}{Student packet, pencil. Teacher models Example 1, then gradual-release into Example 3 (extraneous).}
\end{tabular}
}

\sectionbanner{LAUNCH --- Solving Radical Equations (teacher models on screen)}

\begin{callout}{\IconBook\ ISOLATE--RAISE--CHECK RULES (keep printed on your page)}{calloutgreen}
\begin{enumerate}[leftmargin=2.1em,itemsep=1pt,topsep=3pt]
  \item Isolate the radical on one side of the equation.
  \item Raise BOTH sides to the index power (square for \(\sqrt{\vphantom{x}}\), cube for \(\sqrt[3]{\vphantom{x}}\), etc.).
  \item Solve the resulting polynomial equation.
  \item CHECK each candidate in the ORIGINAL equation --- reject any that don't satisfy it (extraneous).
\end{enumerate}
\end{callout}

\begin{callout}{\IconWarn\ COMMON ERROR}{calloutred}
Do NOT skip the CHECK step --- squaring both sides is not a reversible operation.\\
A candidate solution that makes the radicand negative under an even-index radical is extraneous.\\
Always substitute back into the ORIGINAL equation, not the squared form.
\end{callout}

\sectionbanner{EXPLORE --- Think-Pair-Share (33 min)}
\textit{Work with a partner. Use the Isolate--Raise--Check Rules above for every problem. Move through items in order --- each builds on the last.}

\textbf{Bank-item labels (in order):}
\begin{enumerate}[leftmargin=1.6em,itemsep=2pt,topsep=2pt]
  \item Try It 1 --- Isolate and raise to solve
  \item Try It 3 --- Solve and check for extraneous solutions
  \item Practice \#22 --- Solve radical equation
  \item Practice \#8 --- Solve radical equation
  \item Practice \#15 --- Identify extraneous solutions
  \item Practice \#4 --- Solve and verify
\end{enumerate}

\bankitem{Try It 1 --- Isolate and raise to solve}{%
Solve each radical equation. Check your solution by substituting into the given equation or graphing.

\begin{enumerate}[label=\textbf{\alph*.},leftmargin=1.8em,itemsep=3pt,topsep=4pt]
  \item \(\sqrt{x - 2} + 3 = 5\)
  \item \(\sqrt[3]{x - 1} = 2\)
\end{enumerate}
}{}

\bankitem{Try It 3 --- Solve and check for extraneous solutions}{%
Solve each radical equation. Check for any extraneous solutions.

\begin{enumerate}[label=\textbf{\alph*.},leftmargin=1.8em,itemsep=3pt,topsep=4pt]
  \item \(x = \sqrt{7x + 8}\)
  \item \(x + 2 = \sqrt{x + 2}\)
\end{enumerate}
}{}

\bankitem{Practice \#22 --- Solve radical equation}{%
Solve each radical equation. \((\sqrt{4x}=11)\)
}{}

\bankitem{Practice \#8 --- Solve radical equation}{%
\(\sqrt{8x + 9} = x\)
}{}

\bankitem{Practice \#15 --- Identify extraneous solutions}{%
Generalize Explain how to identify an extraneous solution for an equation containing a radical expression.
}{}

\bankitem{Practice \#4 --- Solve and verify}{%
Error Analysis Nazim said that -3 and 6 are the solutions to \(\sqrt{3x + 18} = x\). What error did Nazim make?
}{}

\sectionbanner{SHARE / SUMMARY (5 min)}

\begin{sentenceframebox}
\textbf{Sentence frames}

Pair share: ``I isolated the radical to get \_\_\_. Squaring both sides gives \_\_\_. My candidate solutions are \_\_\_. When I check in the ORIGINAL equation, \_\_\_ is valid and \_\_\_ is extraneous because \_\_\_.'' 

Whole class: one pair reads their sentence. Class signals thumbs up/sideways.
\end{sentenceframebox}

\sectionbanner{EXIT}

\begin{summaryexitbox}{EXIT TICKET --- Summary Recap}
To solve a radical equation I first \_\_\_.

I raise both sides to power \_\_\_ because \_\_\_.

An extraneous solution happens when \_\_\_.

One biggest thing I learned today: \_\_\_\_\_\_\_\_\_\_\_\_\_\_\_\_\_\_\_\_\_\_
\end{summaryexitbox}

\clearpage

\sectionbanner{OPTIONAL REINFORCEMENT (not collected)}
\textit{If you finish Explore early or want extra practice before Period 2.}

\begin{callout}{\IconSearch\ OPTIONAL \textperiodcentered\ Model \& Discuss}{calloutblue}
Return to the Do Now equation. With your partner:

B) Use Structure. Why can squaring both sides introduce a solution that doesn't work in the original equation?

C) Examine at least two equations you solved in Explore. For each, explain in your own words why checking is a required step and not optional.
\end{callout}

\bankitem{Optional \#1  (Savvas Practice)}{%
Solve each radical equation. \((\sqrt[3]{x}+8=13)\)
}{}

\end{document}
